\documentclass{article}
\usepackage[utf8]{inputenc}
\usepackage[margin=0.1in]{geometry}
\usepackage{listings}
\usepackage{xcolor}
\usepackage{multicol}

% Listings: same lexer/styles as exam-test/main.tex (HTML / CSSEXAM / JSEXAM)
\lstdefinelanguage{CSSEXAM}{%
  morekeywords={%
    color,background,background-color,margin,padding,border,border-width,%
    border-style,border-color,width,height,display,font,font-size,font-style,%
    font-weight,text-align,text-decoration,box-sizing,position,overflow,%
    opacity,max-width,min-width,flex,align-items,justify-content,gap,z-index},%
  sensitive=false,%
  morecomment=[s]{/*}{*/},%
  morestring=[b]',%
}%
\lstdefinelanguage{JSEXAM}{%
  morekeywords={%
    break,case,catch,class,const,continue,debugger,default,delete,do,else,%
    export,extends,false,finally,for,function,if,import,in,instanceof,new,%
    null,return,super,switch,this,throw,true,try,typeof,var,void,while,%
    let,static,async,await,of,yield},%
  sensitive=true,%
  morecomment=[l]{//},%
  morecomment=[s]{/*}{*/},%
  morestring=[b]',%
  morestring=[b]"% 
}%
\lstdefinestyle{htmlexam}{%
  language=HTML,
  basicstyle=\ttfamily\footnotesize,
  keywordstyle=\bfseries\color{blue!55!black},
  commentstyle=\ttfamily\color{black!50},
  stringstyle=\color{green!50!black},
  columns=fullflexible,
  keepspaces=true,
  tabsize=2,
  showstringspaces=false,
  breaklines=true,
  breakatwhitespace=true,
  frame=single,
  rulecolor=\color{black!35},
  backgroundcolor=\color{black!4},
  xleftmargin=0.5em,
  xrightmargin=0.5em,
  aboveskip=4pt,
  belowskip=4pt,
}%
\lstdefinestyle{cssexam}{%
  language=CSSEXAM,
  basicstyle=\ttfamily\footnotesize,
  keywordstyle=\bfseries\color{blue!55!black},
  commentstyle=\ttfamily\color{black!50},
  stringstyle=\color{green!50!black},
  columns=fullflexible,
  keepspaces=true,
  tabsize=2,
  showstringspaces=false,
  breaklines=true,
  breakatwhitespace=true,
  frame=single,
  rulecolor=\color{black!35},
  backgroundcolor=\color{black!4},
  xleftmargin=0.5em,
  xrightmargin=0.5em,
  aboveskip=4pt,
  belowskip=4pt,
}%
\lstdefinestyle{jsexam}{%
  language=JSEXAM,
  basicstyle=\ttfamily\footnotesize,
  keywordstyle=\bfseries\color{blue!55!black},
  commentstyle=\ttfamily\color{black!50},
  stringstyle=\color{green!50!black},
  columns=fullflexible,
  keepspaces=true,
  tabsize=2,
  showstringspaces=false,
  breaklines=true,
  breakatwhitespace=true,
  frame=single,
  rulecolor=\color{black!35},
  backgroundcolor=\color{black!4},
  xleftmargin=0.5em,
  xrightmargin=0.5em,
  aboveskip=4pt,
  belowskip=4pt,
}%

\title{Web Development (HTML, CSS, JS) Cheatsheet}
\author{Software Engineering Technical Reference}
\date{}

\begin{document}

% \maketitle

\begin{multicols}{2}

\section{HTML}

\subsection{Document Structure}
\begin{lstlisting}[style=htmlexam]
<!DOCTYPE html>
<html>
  <head>
    <meta charset="utf-8">
    <title>Page Title</title>
    <!-- Metadata, styles, scripts -->
    <link rel="stylesheet" href="styles.css">
    <script src="app.js" defer></script>
  </head>
  <body>
    <!-- Visible content goes here -->
  </body>
</html>
\end{lstlisting}

\subsection{Elements \& Semantic Markup}
\begin{lstlisting}[style=htmlexam]
<!-- Links & Images (Void element) -->
<a href="page.html" target="_blank">Open</a>
<img src="pic.png" alt="Accessible desc">

<!-- Lists -->
<ul> <!-- Unordered (bullets), ol = Ordered -->
  <li>Item A</li>
  <li>Item B</li>
</ul>

<!-- Semantic blocks vs generic divs -->
<header></header> <nav></nav>
<main></main> <article></article>
<div>Non-semantic container</div>
\end{lstlisting}

\subsection{Tables \& Forms}
\begin{lstlisting}[style=htmlexam]
<table>
  <tr><th>Header 1</th><th>Header 2</th></tr>
  <tr><td>Row 1, C1</td><td>Row 1, C2</td></tr>
</table>

<form>
  <label for="user">Username:</label>
  <input id="user" type="text" name="user">
  <input type="password" name="pass">
  <button type="submit">Login</button>
</form>
\end{lstlisting}

\columnbreak

\section{CSS}

\subsection{Selectors \& Specificity}
\begin{lstlisting}[style=cssexam]
* { box-sizing: border-box; }      /* Universal */
h1 { color: blue; }                /* Element */
.highlight { color: green; }       /* Class */
#main { color: red; }              /* ID (Wins) */
nav a { text-decoration: none; }   /* Descendant */
\end{lstlisting}
Specificity order: Inline \texttt{style=""} $>$ \texttt{\#id} $>$ \texttt{.class} $>$ \texttt{element}.

\subsection{Box Model \& Styling}
\begin{lstlisting}[style=cssexam]
.box {
  /* Box Model: Margin -> Border -> Padding -> Content */
  margin: 16px;             /* Outside border */
  border: 2px solid red;
  padding: 8px;             /* Inside border */
  background-color: yellow; /* Covers content+padding */

  /* Typography */
  text-align: center;       /* Horizontally centers inline content */
  font-style: italic;
  font-size: 14px;
}
\end{lstlisting}

\columnbreak

\section{JavaScript Fundamentals}

\subsection{Variables \& Types}
\begin{lstlisting}[style=jsexam]
// Temporal Dead Zone (let and const)
let count = 0;      // Block-scoped, mutable
const MAX = 10;     // Block-scoped, immutable ref

console.log(myCat); // Throws ReferenceError!
let myCat = "Whiskers";

// Hoisted Zone (Function and var)
var legacy = 1;     // Function-scoped, hoisted (malloc -> execute)
console.log(myDog); // hoisted: undefined (Not a ReferenceError!)
var myDog = "Rex";
console.log(myDog); // Outputs: "Rex"

typeof 42;          // "number" (only 1 numeric type)
typeof "42";        // "string"
\end{lstlisting}

\subsection{Equality \& Falsy Values}
\begin{lstlisting}[style=jsexam]
5 === "5"; // false (Strict: no type coercion)
5 == "5";  // true  (Loose: allows type coercion)

// Falsy values: false, 0, "", null, undefined, NaN
// Truthy values include: [], {}, "0"
\end{lstlisting}

\subsection{Arrays (mutable)}
\begin{lstlisting}[style=jsexam]
const arr = [1, 2, 3];
arr.push(4);                     // Adds to end; [1,2,3,4]
arr.pop();                       // Removes end; returns 4; [1,2,3]
let last = arr[arr.length - 1];  // 3 (zero-indexed)

let part = arr.slice(1, 3);      // [2,3] copy; arr unchanged
arr.includes(2);                 // true (on [1,2,3])
arr.indexOf(9);                  // -1 if missing

// splice(start, deleteCount, ...insert): mutates arr
let cut = arr.splice(1, 1); // [1,2,3]->[1,3]; cut==[2]
arr.splice(1, 0, 99);       // [1,3]->[1,99,3]
\end{lstlisting}

\subsection{Array Methods (pipeline)}
\begin{lstlisting}[style=jsexam]
const items = [
  { name: "A", price: 100, instock: true },  // -> 90
  { name: "B", price: 50,  instock: false }, // out of stock
  { name: "C", price: 200, instock: true },  // -> 180
  { name: "D", price: 10,  instock: true },  // -> 9
];
// filter: keep matches -> new array (can shrink)
// map:    transform each -> new array (same length as input)
// reduce: fold to one value; (accumulator, item) => ...; start
const total = items
  .filter(p => p.instock === true) // [A,C,D]; B removed
  .map(p => p.price * 0.9)         // [90, 180, 9]
  .reduce((sum, x) => sum + x, 0); // 90+180+9 -> 279 

items.find(p => p.name === "C");   // {name:"C",...} or undefined
items.some(p => !p.instock);       // true (B is out of stock)
items.every(p => p.instock);       // false (not all in stock)
\end{lstlisting}

\columnbreak



\subsection{Loops \& Early Exit}
\begin{lstlisting}[style=jsexam]
// Index loop: need i or walk nested structures
for (let i = 0; i < arr.length; i++) {
  const item = arr[i];
}

// for...of: values only (index not required)
for (const item of arr) { /* ... */ }

// Skip to next iteration (guard / filter in-loop)
for (const u of users) {
  if (u.isActive !== true) continue;
}

// Nested loop: each item has its own inner array
for (let i = 0; i < users.length; i++) {
  for (let j = 0; j < users[i].roles.length; j++) {
    const role = users[i].roles[j];
  }
}

if (!arr || arr.length === 0) return null; // empty input
if (typeof raw !== "string") return [];    // wrong type
\end{lstlisting}

\subsection{Objects: Access, Count \& Group}
\begin{lstlisting}[style=jsexam]
const user = { name: "Ada", score: 10, "user-id": 7 };
user.name;              // dot: key is a simple identifier
user["user-id"];        // bracket: key has hyphen/space/symbol
const field = "score";
user[field];            // bracket: key stored in a variable

// Counter object: each key -> how many times seen
const words = ["a", "b", "a", "a", "b"];
const freq = {};
let best = words[0], max = 0;
for (const s of words) {
  if (freq[s] === undefined) freq[s] = 0;
  freq[s]++;             // a:1->2->3, b:1->2
  if (freq[s] > max) {
    max = freq[s]; best = s;
  }
}
// freq == {a:3, b:2}; best == "a"

// Bucket object: each key -> array of values
const groups = {};
if (groups["Dev"] === undefined) groups["Dev"] = [];
groups["Dev"].push("Alice");       // { Dev: ["Alice"] }
groups["Dev"].push("Bob");         // { Dev: ["Alice","Bob"] }
if (groups["Admin"] === undefined) groups["Admin"] = [];
groups["Admin"].push("Carol");     // adds second role key

for (const k of Object.keys(groups)) { 
  /* k of ["Dev", "Admin"] */
}
for (const v of Object.values(groups)) {
  /* v of [["Alice","Bob"], ["Carol"]] */
}
for (const [role, names] of Object.entries(groups)) {
  /* role of ["Dev", "Admin"], names of ["Alice","Bob"] or ["Carol"] */
}
\end{lstlisting}

\columnbreak

\subsection{Strings \& JSON}
\begin{lstlisting}[style=jsexam]
const raw = "  A, B ,, c "; 
raw.split(",");             // ["  A"," B ",""," c "]

"  hi ".trim();             // "hi"
"API".toLowerCase();        // "api"

`Score: ${user.score}!`;  // "Score: 10!" (backticks interpolation)
"Score: ${user.score}!";  // "Score: ${user.score}!"

// split -> trim each -> skip "" -> lowercase -> collect
const parts = raw.split(",");
const tokens = [];
for (let i = 0; i < parts.length; i++) {
  const t = parts[i].trim(); 
  if (t !== "") tokens.push(t.toLowerCase());
}
// tokens == ["a", "b", "c"]

// JSON = text format; these convert object <-> string
const user = { name: "Ada", score: 10, "user-id": 7 };
const str = JSON.stringify(user);   // object -> String (text)
// str looks like: {"name":"Ada","score":10,"user-id":7}
const obj = JSON.parse(str);        // String (text) -> object
\end{lstlisting}




\section{DOM Manipulation \& Events}

\subsection{Selecting \& Modifying Elements}
\begin{lstlisting}[style=jsexam]
// <button id="go"> <p class="active"> <input id="user">
// <img id="pic" src="old.png"> <div data-user-id="7">

const btn = document.getElementById("go");     // by id
const el = document.querySelector(".active");  // CSS selector
const inputEl = document.getElementById("user");
const pic = document.getElementById("pic");

el.textContent = "Done";        // plain text; shows <b> as chars
el.innerHTML = "<b>x</b>";      // parses HTML -> bold x
const val = inputEl.value;      // text user typed in <input>

pic.setAttribute("src", "a.png"); // change attribute (was old.png)
el.classList.add("highlight");    // add CSS class from stylesheet
// el.classList.remove("highlight"); el.classList.toggle("on");

const box = document.querySelector("[data-user-id]");
let id = box.dataset.userId;  // data-user-id="7" -> "7"
\end{lstlisting}

\subsection{Creating Elements \& Events}
\begin{lstlisting}[style=jsexam]
// <ol id="list"><li>Old</li></ol>

const list = document.getElementById("list");
const li = document.createElement("li"); // built, not on page yet
li.textContent = "New item";
list.appendChild(li);          // attach child -> shows in list
li.remove();                   // detach from parent

// <a id="go" href="/home">Home</a>  <- click normally opens /home
const link = document.getElementById("go");
link.addEventListener("click", (event) => {
  event.preventDefault();      // block going to /home
  console.log(event.target);   // <a id="go"> (what user clicked)
  // your code runs instead (e.g. update page with JS)
});
\end{lstlisting}


\columnbreak

\section{Advanced JavaScript}

\subsection{Classes \& Error Handling}
\begin{lstlisting}[style=jsexam]
class Widget {
  constructor(name) { 
    this.name = name; 
  } 
  label() { return this.name; }  // method on instance
}
const w = new Widget("clock");   // w.name == "clock"
w.label();                       // "clock"

try {
  JSON.parse("not json");        // throws if invalid
} catch (e) {
  console.log(e.message);        // handle error, program continues
}
\end{lstlisting}

\subsection{Rest \& Spread Syntax}
\begin{lstlisting}[style=jsexam]
// REST (... in params): csv -> array
function sum(...nums) {
  return nums.reduce((a, b) => a + b, 0);
}
sum(3, 9, 4);                // nums == [3, 9, 4] -> 16

// SPREAD (... on array): array -> csv
const arr = [1, 2];
const merged = [...arr, 3];  // unpack arr -> [1, 2, 3]
// [arr, 3] would be [[1,2], 3] (nested, not what you want)

const scores = [88, 92, 75]; // already an array (from data/API)
Math.max(scores);            // NaN (function max(...nums))
Math.max(...scores);         // 92
\end{lstlisting}

\subsection{Promises \& Async / Await}
\begin{lstlisting}[style=jsexam]
setTimeout(() => console.log("Later"), 1000);
// code keeps running now; "Later" prints after 1 second

// fetch: network request (async = finishes later)
async function loadData(url) {
  // step 1: Response (status, headers)
  const res = await fetch(url);
  // 404: fetch ok, res.ok false
  if (!res.ok) throw new Error(`HTTP ${res.status}`);
  // step 2: read JSON body -> object/array
  const data = await res.json();
  // e.g. [{id:1}, {id:2}]
  return data;
}

async function main() {
  const data = await loadData("/api/items"); // pause until done
}

// Need TWO URLs? Compare one-by-one vs together:
async function loadBoth() {
  const u1 = "/api/users", u2 = "/api/scores";

  // SLOW: await fetch(u1); then await fetch(u2);
  //        |--- 200ms ---|   |--- 200ms ---|  ~400ms total

  const results = await Promise.all([
    fetch(u1),   // \
    fetch(u2),   //  > start together; wait for both
  ]);            //     ~200ms total
  const resUsers = results[0];   // same order as above
  const resScores = results[1];
  const users = await resUsers.json();
  const scores = await resScores.json();
}
\end{lstlisting}

\end{multicols}

\end{document}
